% ===========================================================
%  CHAPTER I — GENERAL INTRODUCTION
% ===========================================================
\chapter{General Introduction}
\label{chap:intro}

% -----------------------------------------------------------
\section{Context and Motivation}

\subsection{A Fundamental Scientific Question}

Human movement, in all its richness and complexity, results from a remarkable
orchestration between the central nervous system~(CNS), the musculoskeletal
structures, and the sensory organs. Understanding the mechanisms that govern this
motor control is a major scientific challenge, with considerable implications for
rehabilitation medicine~\parencite{Kwakkel2003}, bio-inspired robotics, the design of
myoelectric prostheses, and neurocomputational modeling.

Musculoskeletal models provide a formal framework for simulating the biomechanics of
movement. Platforms such as OpenSim~\parencite{Delp2007} make it possible to
faithfully reconstruct the anatomy and dynamics of the limbs. However, the optimal
control of such models remains a fundamentally difficult problem --- often referred to
as the ``inverse problem of movement'' --- because of (i)~motor
redundancy~\parencite{Zajac1989}, (ii)~the strong nonlinearity of contractile
elements, and (iii)~the high dimensionality of the state and action spaces.

At the same time, the emergence of \textit{deep reinforcement learning}~(\DRL{}) has
opened new perspectives for the automated learning of control strategies in complex
environments~\parencite{Mnih2015, Sutton2018}. Pioneering work has demonstrated the
ability of \DRL{} algorithms to master continuous-control tasks of remarkable
sophistication~\parencite{Lillicrap2016, Schulman2017, Haarnoja2018}. More recently,
Meta~AI's MyoSuite project~\parencite{Caggiano2022}, the
MotorNet platform~\parencite{Codol2023}, and the \textit{Learning to
Run} challenge~\parencite{Kidzinski2018} have established the feasibility of
combining musculoskeletal models and \DRL{} to reproduce biologically plausible human
motor strategies.

These converging developments provide the framework within which the present work is
situated.

\subsection{Epidemiological and Clinical Challenges}
\label{sec:epidemio}

Beyond its fundamental scientific interest, the societal stake is considerable.
Stroke is the world's second leading cause of death and a leading cause of
acquired disability in adults~\parencite{Feigin2022}.

An earlier version of this section also reported a national incidence figure for
Tunisia. That figure was attributed to a source whose identifier resolves to an
unrelated article, and no substitute could be verified; both the statistic and
its citation have therefore been withdrawn rather than re-attributed. The
argument does not depend on it --- the global burden reported
by~\textcite{Feigin2022} is sufficient motivation, and it is independently
verifiable.

Among stroke survivors, more than $60$~\% still present residual upper-limb motor
impairment six months after the event~\parencite{Kwakkel2003}. Intensive,
personalized rehabilitation, guided by robotic devices and assisted therapies,
significantly improves recovery~\parencite{Lo2010}. Nevertheless, access to these
technologies remains limited by their cost and by the shortage of trained therapists,
particularly in developing countries. A musculoskeletal simulation environment capable
of exploring rehabilitation strategies \textit{in silico} before clinical application
therefore represents a major scientific and humanitarian opportunity.

% -----------------------------------------------------------
\section{Problem Statement}

The central problem addressed in this thesis can be formulated as follows:

\begin{quote}
\itshape For a reaching movement of a nine-muscle upper limb, does a controller
trained by deep reinforcement learning distribute force across the individual
muscles in the same way as classical static optimisation --- recruiting the same
muscles in the same proportions, and at the same total cost --- and under what
conditions on the training objective does that agreement hold?
\end{quote}

This formulation replaces the one carried by earlier drafts, which asked whether
a trained policy could \emph{replace the iterative Newton--Raphson solver at
inference time, trading a one-off training cost against virtually zero per-step
computation}. That question presupposed its own answer, and the presupposition
did not survive measurement: \cref{sec:latencycaveat} finds that a competently
implemented classical load-sharing solver runs in \SI{23.9}{\micro\second}
against the network's \SI{288.8}{\micro\second}, so the network is roughly twelve
times \emph{slower}, not faster. The speed argument is withdrawn in full, and the
problem statement is rewritten around the question the evidence actually answers.

This problem can be broken down into five research questions:
\begin{enumerate}[label=\textbf{Q\arabic*.}, leftmargin=2.2em]
  \item What level of physiological fidelity is necessary and sufficient for a
        musculoskeletal model dedicated to \DRL{}?
  \item Which reward-function formulations encourage biologically plausible motor
        strategies while preserving learning efficiency?
  \item Which \DRL{} algorithms provide the best compromise between convergence,
        robustness, and energy efficiency on muscular state and action spaces?
  \item Do the trained agents exhibit behaviours reported in humans
        (co-contraction, energy economy, muscle synergies)?
  \item Does Domain Randomization improve the robustness of musculoskeletal policies,
        and to what extent does it make transfer to real patients with varying
        anatomical parameters conceivable?
  \item[\textbf{Q6.}] How does the inference cost of a trained policy compare with
        that of the classical calculation it would displace, measured
        like-for-like on the same problem?
\end{enumerate}

Two of these were narrowed after measurement. \textbf{Q4} originally also promised
an analysis of \textit{jerk} minimisation and of triphasic activation profiles;
neither analysis was performed, and the claim has been removed rather than left
outstanding. \textbf{Q6} originally asked what \emph{computational gain} a policy
delivers, which assumed there was one; it is restated here as a neutral
comparison, and \cref{sec:latencycaveat} reports that the comparison favours the
classical solver.

% -----------------------------------------------------------
\section{Objectives}

This thesis pursues the following objectives:

\begin{enumerate}
  \item Implement an articulated multibody dynamic model of an upper limb in
        \MuJoCo{}~\parencite{Todorov2012} --- four actuated degrees of freedom
        driven by nine muscles --- specifying a Hill-type
        model~\parencite{Hill1938} with CE, SEE, PEE, pennation-angle dynamics
        and numerical integration of fibre--tendon equilibrium.
  \item Validate the physical model by comparing passive mechanical properties and
        FL/FV curves with data from the literature~\parencite{Holzbaur2005, Winter2009, Gordon1966}.
  \item Develop a modular Gymnasium environment with Domain Randomization and
        reproducible perturbations.
  \item Define reference controllers (random policy, impedance PD controller) in
        order to establish a meaningful evaluation baseline.
  \item Compare four \DRL{} algorithms (PPO, SAC, TD3, DDPG) under a common
        protocol, with hyperparameters searched using
        Optuna~\parencite{Akiba2019}. The protocol originally specified
        $10$~seeds at $2{\times}10^6$~steps; \cref{sec:compute} explains why the
        available compute did not permit it, and the comparison as executed uses
        three seeds at $10^5$~steps, with the principal configuration replicated
        across five seeds at $3{\times}10^5$. No significance testing is claimed
        at those sample sizes.
  \item Analyze and quantify emerging motor strategies: co-contraction
        (CCI~\parencite{Falconer1985}), metabolic cost~\parencite{Bhargava2004},
        robustness to perturbations, temporal activation profiles, and muscle
        synergies obtained through non-negative matrix factorization~\parencite{dAvella2003}.
  \item Discuss the implications for post-stroke rehabilitation and biomedical
        robotics, including the ethical considerations specific to these
        applications.
  \item Benchmark the computational cost of the Newton--Raphson muscle
        equilibrium solver (per-step wall-clock time for all 9 muscles) against
        the inference latency of the trained RL policy, quantifying the
        trade-off between physics-based exactness and real-time deployability on
        embedded rehabilitation hardware.
\end{enumerate}

% -----------------------------------------------------------
\section{Original Contributions}
\label{sec:contributions}

Compared with similar studies, this thesis makes the following contributions:

\begin{enumerate}[label=\textbf{C\arabic*.}, leftmargin=2.2em]
  \item \textbf{A fully specified Hill-type model in \MuJoCo{}}: explicit
        treatment of the SEE and PEE equations and of pennation-angle dynamics,
        with numerical solution of fibre--tendon equilibrium by the
        Newton--Raphson method, each component justified mathematically and
        provided as reference code (Appendix~\ref{app:newton}), where comparable
        studies commonly adopt \MuJoCo{}'s internal simplified muscle model
        without a stated validation protocol.

        One qualification belongs with this claim rather than in a later
        limitations section. Probing every muscle at zero activation across
        sampled trajectories returns a passive force of exactly zero in all
        states: no muscle in this model exceeds $0.67$ of its declared length
        range, and the parallel element only engages near the top of that range.
        The specification is complete; the \emph{exercised} model is the
        contractile and series-elastic elements, with the force--length curve
        used on its ascending limb and plateau only. ``Complete Hill model''
        would therefore overstate what is active, and the weaker wording is used
        throughout.
  \item \textbf{A quantitative validation protocol}: comparison of FL and FV
        curves and passive joint torques against reference anthropometric
        data~\parencite{Holzbaur2005, Gordon1966}, with NRMSE reporting.
  \item \textbf{A per-muscle comparison between \DRL{} and classical static
        optimisation} (\cref{ch:forces}). Both methods are posed the identical
        load-sharing problem on the identical trajectory, with the constraint
        residual verified to machine precision, and agreement is reported muscle
        by muscle. This is the contribution around which the present work is
        organised, and it addresses the research question directly rather than
        through task performance.
  \item \textbf{Identification of the discount factor as the decisive
        hyperparameter for this task} (\cref{sec:ablation2}). A one-at-a-time
        ablation over five searched parameters, three seeds each, finds that
        $\gamma$ alone recovers the entire default-to-tuned improvement
        ($+106$\,\%) while the other four recover nothing. The mechanism is
        checkable and quantitative: $\gamma$ sets an effective horizon of about
        $1/(1-\gamma)$ steps, and the selected value matches the duration of the
        reach itself where the library default exceeds it fourfold.

        This contribution replaces one claimed by earlier drafts --- that the
        \SAC{} target entropy of $-\dim(\mathcal{A})$ imposed the error plateau.
        That hypothesis was mechanistically plausible and supported by the
        clustering of the best search trials, and it was wrong: the same ablation
        shows the target entropy recovers $-8$\,\% of the gap, i.e.\ nothing
        (\cref{sec:ablation}). It is reported in full in this thesis as a
        negative result rather than removed.
  \item \textbf{Evidence that configuration dominates algorithm selection}
        (\cref{sec:fairbench}): with every algorithm given a correctly matched
        discount factor, the interval between the best and second-best off-policy
        method is \SI{0.060}{\meter}, while correcting $\gamma$ alone within a
        single algorithm is worth \SI{0.165}{\meter} --- roughly three times the
        gap between algorithms.
  \item \textbf{A documented account of six substantive defects}
        (\cref{sec:defects}), four of which produced plausible-looking but
        incorrect results, together with the measurements that exposed each.
  \item \textbf{A like-for-like inference-cost comparison, and the retraction it
        forced} (\cref{sec:latencycaveat}). An initial measurement gave the policy
        a $7.3\times$ advantage over a Newton--Raphson solve. Re-examined, that
        comparison timed the wrong computation --- the muscle-model solve inside
        forward simulation rather than the load-sharing problem the thesis
        actually poses --- and both sides were unoptimised. Measured
        like-for-like on the load-sharing problem, the classical route takes
        \SI{23.9}{\micro\second} against the network's \SI{288.8}{\micro\second}.
        The speed advantage is withdrawn; what survives is the weaker property
        that the network's cost is fixed and branch-free where the solver's
        depends on convergence.
\end{enumerate}

% -----------------------------------------------------------
\section{Thesis Organization}

This document is organized as follows:
\begin{description}[leftmargin=3.5em, style=nextline]
  \item[\textbf{Chapter~II}] presents the state of the art in musculoskeletal
        modeling, computational theories of motor control, and \DRL{} for
        continuous-action spaces.
  \item[\textbf{Chapter~III}] describes the design of the complete
        musculoskeletal model, its validation, and the simulation environment.
  \item[\textbf{Chapter~IV}] presents the theoretical framework of the \DRL{}
        algorithms used, from the policy-gradient theorem to actor--critic
        architectures.
  \item[\textbf{Chapter~V}] details the implementation and experimental
        methodology, including baselines, the statistical protocol, and
        reproducibility measures.
  \item[\textbf{Chapter~VI}] analyzes and discusses the experimental results,
        from convergence to emerging motor strategies and synergy analysis.
  \item[\textbf{Chapter~VII}] interprets the results, confronts them with the
        literature, addresses ethical considerations, and identifies the
        limitations of the work.
  \item[\textbf{Conclusion}] summarizes the contributions and outlines future
        directions.
\end{description}
